\section*{Hoofdstuk \ref{ch:g10:reaction-progress-table} --- De voortgangstabel}

\begin{solution}{exo:g10:reaction-progress-table:1}
\ce{H2}: $4.0 - 2x$; \ce{O2}: $3.0 - x$; \ce{H2O}: $2x$. Waterstof is op
bij $x = 2.0$, zuurstof bij $x = 3.0$: $x_{\max} = \qty{2.0}{mol}$,
waterstof is beperkend. \omterm{def:g10:reaction-progress-table:system}{Eindtoestand}: \qty{0}{mol} \ce{H2},
\qty{1.0}{mol} \ce{O2}, \qty{4.0}{mol} \ce{H2O}.
\end{solution}

\begin{solution}{exo:g10:reaction-progress-table:2}
Koolstof zou op zijn bij $x = 3.0$, zuurstof bij $x = 2.0$:
$x_{\max} = \qty{2.0}{mol}$, zuurstof is beperkend. \omterm{def:g10:reaction-progress-table:system}{Eindtoestand}:
\qty{1.0}{mol} koolstof, geen zuurstof, \qty{2.0}{mol} koolstofdioxide.
\end{solution}

\begin{solution}{exo:g10:reaction-progress-table:3}
$x_{\max} = \qty{0.30}{mol}$ (zwavel beperkend). \omterm{def:g10:reaction-progress-table:system}{Eindtoestand}:
\qty{0.20}{mol} ijzer, geen zwavel, \qty{0.30}{mol} ijzersulfide.
\end{solution}

\begin{solution}{exo:g10:reaction-progress-table:4}
Stikstof is op bij $x = 1.0$, waterstof bij $x = 2.4/3 = 0.80$:
$x_{\max} = \qty{0.80}{mol}$. \omterm{def:g10:reaction-progress-table:system}{Eindtoestand}: \qty{0.20}{mol} \ce{N2}, geen
\ce{H2}, $2 \times 0.80 = \qty{1.6}{mol}$ \ce{NH3}.
\end{solution}

\begin{solution}{exo:g10:reaction-progress-table:5}
(b): $4/2 = 2/1$. In (a) is koolstofmono-oxide beperkend, in (c) ook.
\end{solution}

\begin{solution}{exo:g10:reaction-progress-table:6}
$n(\ce{CH4}) = 32 / 16.0 = \qty{2.0}{mol}$, dus $x_{\max} = \qty{2.0}{mol}$
voor een \omterm{def:g10:reaction-progress-table:stoichiometric}{stoichiometrisch mengsel}. Zuurstof: $2 \times 2.0 = \qty{4.0}{mol}$,
dat is $4.0 \times 32.0 = \qty{128}{g}$. Water:
$\qty{4.0}{mol} \times \qty{18.0}{g/mol} = \qty{72}{g}$.
\end{solution}

\begin{solution}{exo:g10:reaction-progress-table:7}
$n(\ce{C3H8}) = 1000 / 44.0 = \qty{22.7}{mol}$; koolstofdioxide
$3 \times 22.7 = \qty{68.2}{mol}$; $V = 68.2 \times 24.1 =
\qty{1.64e3}{L}$, ongeveer \qty{1.6}{m^3}.
\end{solution}

\begin{solution}{exo:g10:reaction-progress-table:8}
Bij $x = \qty{1.0}{mol}$: \qty{1.0}{mol} \ce{CH4}, \qty{1.0}{mol}
\ce{O2}, \qty{1.0}{mol} \ce{CO2}, \qty{2.0}{mol} \ce{H2O}. De zuurstof
is op bij $x = \qty{1.5}{mol}$: daar stopt de reactie, dus grotere
voortgangen worden nooit bereikt.
\end{solution}

\begin{solution}{exo:g10:reaction-progress-table:9}
$n_0(\ce{Zn}) = 1.30 / 65.4 = \qty{0.0199}{mol}$;
$n_0(\ce{Cu^{2+}}) = 0.10 \times 0.1000 = \qty{0.0100}{mol}$. De
koperionen zijn beperkend: $x_{\max} = \qty{0.0100}{mol}$. Gevormd koper:
$0.0100 \times 63.5 = \qty{0.635}{g}$; overgebleven zink:
$(0.0199 - 0.0100) \times 65.4 = \qty{0.65}{g}$. Er zijn geen
koperionen meer over: de blauwe kleur is verdwenen.
\end{solution}

\begin{solution}{exo:g10:reaction-progress-table:10}
$n_0(\ce{Mg}) = 0.48 / 24.3 = \qty{0.0198}{mol}$, dat op zou zijn bij
$x = 0.0198$; $n_0(\ce{H+}) = \qty{0.020}{mol}$, dat op zou zijn bij
$x = 0.010$. Het zuur is beperkend: $x_{\max} = \qty{0.010}{mol}$, dus
\qty{0.010}{mol} waterstof, $0.010 \times 24.1 = \qty{0.24}{L}$.
\end{solution}

\begin{solution}{exo:g10:reaction-progress-table:11}
Een \omterm{def:g10:reaction-progress-table:stoichiometric}{stoichiometrisch mengsel} met \qty{0.30}{mol} zwavel heeft
\qty{0.30}{mol} ijzer nodig; er is \qty{0.20}{mol} meer, dat is
$0.20 / 0.30 \approx \qty{67}{\percent}$ overmaat.
\end{solution}

\begin{solution}{exo:g10:reaction-progress-table:12}
Gelijke hoeveelheden: $m(\ce{Fe})/55.8 = m(\ce{S})/32.1$, met
$m(\ce{Fe}) + m(\ce{S}) = \qty{10.0}{g}$. Dus
$m(\ce{Fe}) = 10.0 \times 55.8 / (55.8 + 32.1) = \qty{6.35}{g}$ en
$m(\ce{S}) = \qty{3.65}{g}$.
\end{solution}

\begin{solution}{exo:g10:reaction-progress-table:13}
\begin{enumerate}
  \item $n_0(\ce{CaCO3}) = 10.0 / 100.1 = \qty{0.0999}{mol}$, de enige
        \omterm{def:g7:chemical-reactions:reactant}{beginstof}: $x_{\max} = \qty{0.0999}{mol}$; dat geeft
        \qty{0.0999}{mol} \ce{CaO} en \ce{CO2}, dus
        $0.0999 \times 24.1 = \qty{2.41}{L}$ koolstofdioxide.
  \item $n_0(\ce{CaO}) = \qty{0.0999}{mol}$, $n_0(\ce{H2O}) = 1.8 / 18.0
        = \qty{0.100}{mol}$: de ongebluste kalk is (net) beperkend; het
        \omterm{def:g6:pure-substances-and-mixtures:mixture}{mengsel} is bijna stoichiometrisch. \ce{Ca(OH)2}:
        $M = 40.1 + 2 \times (16.0 + 1.0) = \qty{74.1}{g/mol}$, massa
        $0.0999 \times 74.1 = \qty{7.40}{g}$.
\end{enumerate}
\end{solution}

\begin{solution}{exo:g10:reaction-progress-table:14}
$n_0(\ce{H+}) = \qty{0.050}{mol}$ in elke kolf; dat zou op zijn bij
$x = \qty{0.025}{mol}$. Magnesium: 0.15, 0.30, 0.45, 0.60, 0.75,
\qty{0.90}{g} geven 0.0062, 0.0123, 0.0185, 0.0247, 0.0309,
\qty{0.0370}{mol}. In de eerste vier kolven is magnesium beperkend en is
de hoeveelheid waterstof $n(\ce{Mg}) \times 24.1$: 0.15, 0.30, 0.45 en
\qty{0.60}{L}. In de laatste twee is het zuur beperkend: $0.025 \times 24.1
= \qty{0.60}{L}$. De grafiek is een rechte lijn door de oorsprong tot
ongeveer \qty{0.61}{g} magnesium, en daarna een horizontaal plateau op
\qty{0.60}{L}.
\end{solution}

\begin{solution}{exo:g10:reaction-progress-table:15}
$n_0(\text{ethanol}) = 4.6 / 46.0 = \qty{0.100}{mol}$ (zou op zijn bij
$x = 0.100$); $n_0(\ce{O2}) = 4.8 / 24.1 = \qty{0.199}{mol}$ (zou op zijn
bij $x = 0.199/3 = 0.0664$). Zuurstof is beperkend,
$x_{\max} = \qty{0.0664}{mol}$. Overgebleven ethanol: $0.100 - 0.0664 =
\qty{0.0336}{mol}$, dat is $0.0336 \times 46.0 = \qty{1.5}{g}$.
Koolstofdioxide: $2 \times 0.0664 = \qty{0.133}{mol}$, dat is
$0.133 \times 24.1 = \qty{3.2}{L}$.
\end{solution}

\begin{solution}{pb:g10:reaction-progress-table:1}
\textbf{1.} Links: 2 Na, 6 N. Rechts: 2 Na, $3 \times 2 = 6$ N. Klopt.

\textbf{2.} Links: 10 Na, 2 K, 2 N, 6 O. Rechts: K: 2; Na: $5 \times 2 =
10$; O: $1 + 5 = 6$; N: 2. Klopt.

\textbf{3.} Het hoofdstuk over het \omterm{def:g9:periodic-table-first-look:periodic-table}{periodiek systeem}
(\cref{ch:g9:periodic-table-first-look}): natrium reageert heftig met
water, met een bijtende oplossing en waterstof als gevolg. Na de botsing
komt de zak in contact met de huid, de ogen en de vochtige \omterm{def:g4:air-a-mixture-of-gases:air}{lucht} van de
adem.

\textbf{4.} $M(\ce{NaN3}) = 23.0 + 3 \times 14.0 = \qty{65.0}{g/mol}$;
$M(\ce{KNO3}) = 39.1 + 14.0 + 3 \times 16.0 = \qty{101.1}{g/mol}$.

\textbf{5.} \ce{NaN3}: $1.00 - 2x$; \ce{Na}: $2x$; \ce{N2}: $3x$.

\textbf{6.} $1.00 - 2x = 0$ bij $x_{\max} = \qty{0.500}{mol}$;
natriumazide, de enige \omterm{def:g7:chemical-reactions:reactant}{beginstof}, is beperkend.

\textbf{7.} \ce{Na}: $2 \times 0.500 = \qty{1.00}{mol}$; \ce{N2}:
$3 \times 0.500 = \qty{1.50}{mol}$.

\textbf{8.} $1.50 \times 24.1 = \qty{36.2}{L}$.

\textbf{9.} \ce{Na}: $1.00 - 10x$; \ce{KNO3}: $n - 2x$; \ce{K2O}: $x$;
\ce{Na2O}: $5x$; \ce{N2}: $x$.

\textbf{10.} $1.00/10 = n/2$, dus $n = \qty{0.200}{mol}$, dat is
$0.200 \times 101.1 = \qty{20.2}{g}$.

\textbf{11.} $x_{\max} = \qty{0.100}{mol}$: geen natrium, geen
kaliumnitraat, \qty{0.100}{mol} \ce{K2O}, \qty{0.500}{mol} \ce{Na2O},
\qty{0.100}{mol} \ce{N2}.

\textbf{12.} Natrium zou op zijn bij $x = 0.100$, kaliumnitraat bij
$x = 0.150/2 = 0.075$: het nitraat is beperkend, $x_{\max} =
\qty{0.075}{mol}$, en er zou $1.00 - 10 \times 0.075 = \qty{0.25}{mol}$
natriummetaal in de zak achterblijven, klaar om met vocht te reageren.

\textbf{13.} $1.50 + 0.100 = \qty{1.60}{mol}$ stikstof per \omterm{def:g10:the-mole:amount}{mol}
natriumazide.

\textbf{14.} $n = 60 / 24.1 = \qty{2.49}{mol}$.

\textbf{15.} $2.49 / 1.60 = \qty{1.56}{mol}$ natriumazide.

\textbf{16.} $1.56 \times 65.0 = \qty{101}{g}$.

\textbf{17.} $0.200 \times 1.56 = \qty{0.311}{mol}$, dat is
$0.311 \times 101.1 = \qty{31.5}{g}$ kaliumnitraat.

\textbf{18.} $2.49 / 1.50 = \qty{1.66}{mol}$, dat is
$1.66 \times 65.0 = \qty{108}{g}$: de tweede reactie bespaart ongeveer
\qty{7}{g} natriumazide.

\textbf{19.} Een \omterm{def:g10:reaction-progress-table:limiting}{aflopende reactie} stopt pas als de \omterm{def:g10:reaction-progress-table:limiting}{beperkende beginstof}
op is; hier is natriumazide de enige \omterm{def:g7:chemical-reactions:reactant}{beginstof}, dus in de \omterm{def:g10:reaction-progress-table:system}{eindtoestand} is
de hoeveelheid ervan $1.00 - 2x_{\max} = 0$. Er blijft niets giftigs in
de zak achter.

\textbf{20.} Ongeveer \qty{101}{g} natriumazide vult een airbag van
\qty{60}{L}.
\end{solution}
